\chapter{Kecerdasan Umum, Superintelligence, dan Keselamatan AI}
\label{ch:keselamatan}

Tidak ada satu mata kuliah pun yang saya ikuti bernama keselamatan AI. Tetapi pertanyaan-pertanyaannya muncul di banyak kuliah:
di AI and law ketika membahas model dengan risiko sistemik, di robopsychology ketika membahas kepercayaan, di reinforcement
learning ketika agen menemukan cara curang untuk mendapat ganjaran, dan di explainable AI ketika kita sadar betapa sedikit yang kita
ketahui tentang isi model besar. Di luar ruang kuliah, perdebatannya jauh lebih ramai. Istilah seperti AGI dan superintelligence
kini muncul di berita, di laporan pemerintah, dan di pernyataan para pendiri laboratorium AI. Bab ini mencoba memetakan perdebatan
itu dengan jujur: apa yang dimaksud dengan istilah-istilahnya, mengapa sebagian orang khawatir, mengapa sebagian lain
skeptis, apa yang sedang dikerjakan secara konkret, dan bagaimana dunia mencoba mengaturnya. Isinya bersandar pada bacaan, bukan
pada satu kuliah, dan bagian yang masih diperdebatkan ditulis sebagai perdebatan.

\section{Istilah-istilah}

Hampir semua sistem di buku ini adalah AI \emph{sempit}: baik pada satu jenis tugas, seperti mengenali tumor, memainkan Go,
atau meramal cuaca, dan tidak berguna di luar tugas itu. \istilah{Artificial general intelligence} (AGI) merujuk pada sistem yang
mampu mengerjakan berbagai tugas kognitif setara manusia. Definisinya sendiri diperdebatkan. \citet{morris2023} menganalisis
definisi-definisi yang ada dan mengusulkan tingkatan AGI menurut dua sumbu: seberapa \emph{baik} sebuah sistem dibanding manusia
pada tugas tertentu, dan seberapa \emph{luas} rentang tugas yang dikuasainya, ditambah sumbu ketiga tentang seberapa otonom ia
dibiarkan bertindak. Kerangka seperti ini membantu, karena pertanyaan ``apakah ini sudah AGI?'' berubah menjadi pertanyaan yang lebih
bisa diukur: seberapa baik, pada tugas apa saja, dengan otonomi seberapa besar.

\istilah{Superintelligence}, dalam rumusan \citet{bostrom2014}, adalah kecerdasan yang jauh melampaui kemampuan kognitif manusia di
hampir semua bidang. Bostrom membahas beberapa jalan menuju ke sana dan mengapa transisinya mungkin terjadi cepat: sistem yang cukup
pintar untuk memperbaiki dirinya sendiri bisa mempercepat kemajuannya sendiri. Apakah dan kapan ini terjadi tidak diketahui siapa
pun, dan ramalan para ahli tersebar lebar.

\section{Mengapa sebagian orang khawatir}

Inti kekhawatiran bisa dirumuskan tanpa fiksi ilmiah. Sebuah sistem yang mengoptimalkan tujuan akan mengejar tujuan itu, bukan
tujuan yang \emph{dimaksudkan} perancangnya. Di \cref{ch:rl} dan di bagian evaluasi \cref{ch:hayati} kita sudah melihat versi kecilnya:
agen yang menemukan celah dalam fungsi ganjaran, dan model generatif yang menemukan celah dalam model penilai. Gejala ini disebut
\istilah{reward hacking} atau \emph{specification gaming}, dan ia adalah hukum Goodhart dalam bentuk algoritmik: ukuran yang dijadikan
target berhenti menjadi ukuran yang baik. \citet{amodei2016} merumuskan lima masalah praktis yang muncul dari sini. Menghindari efek
samping yang tidak diinginkan, menghindari reward hacking, pengawasan yang tetap bisa dilakukan ketika tujuan terlalu mahal untuk
dinilai setiap saat, eksplorasi yang aman, dan tetap berperilaku baik ketika distribusi data bergeser.

\citet{bostrom2014} menambahkan dua tesis. Tesis ortogonalitas menyatakan bahwa tingkat kecerdasan dan isi tujuan bisa bebas satu
sama lain: sistem yang cerdas tidak otomatis punya tujuan yang baik. Tesis konvergensi instrumental menyatakan bahwa untuk
hampir tujuan apa pun, beberapa subtujuan selalu berguna, seperti mempertahankan diri, mengumpulkan sumber daya, dan menolak tujuannya
diubah, karena semuanya membantu mencapai tujuan utama. Gabungan keduanya mengarah pada \emph{masalah kendali}: bagaimana memastikan
sistem yang jauh lebih cerdas tetap bisa dikoreksi.

\citet{russell2019} menawarkan arah jalan keluar. Alih-alih memberi mesin tujuan yang tetap dan dianggap benar, mesin dirancang untuk
memaksimalkan terwujudnya preferensi manusia, sambil tetap \emph{tidak yakin} apa preferensi itu, dan mempelajarinya dari perilaku
manusia. Mesin yang tidak yakin tentang tujuannya punya alasan untuk mengizinkan dirinya dimatikan, karena manusia yang mematikannya
mungkin tahu sesuatu yang belum ia ketahui. Pelatihan model bahasa dengan preferensi manusia di \cref{ch:terbaru} adalah langkah
kecil, dan masih tidak sempurna, ke arah itu.

Kekhawatiran ini tidak hanya datang dari filsuf. Pada 30 Mei 2023, Center for AI Safety menerbitkan pernyataan satu kalimat yang
ditandatangani ratusan peneliti AI dan pimpinan laboratorium besar: mengurangi risiko kepunahan akibat AI harus menjadi prioritas
global, sejajar dengan risiko berskala masyarakat lain seperti pandemi dan perang nuklir. \citet{hendrycks2023} menyusun sumber risiko
katastrofik ke dalam empat kelompok: penyalahgunaan yang disengaja, perlombaan antar pengembang atau negara yang mendorong peluncuran
sistem yang belum aman, risiko organisasi seperti kecelakaan dan budaya keselamatan yang lemah, dan sistem yang lepas kendali.

\section{Mengapa sebagian lain skeptis}

Kritik terhadap narasi ini sama seriusnya. \citet{mitchell2021} mengingatkan bahwa sejarah AI penuh siklus antara ramalan yang terlalu
optimistis dan kekecewaan. Ia menunjuk beberapa kekeliruan yang berulang: menganggap kemajuan pada tugas sempit sebagai langkah
pasti menuju kecerdasan umum, mengira bahwa yang sulit bagi manusia juga sulit bagi mesin dan sebaliknya, terkecoh oleh nama-nama
yang memberi kesan pemahaman, seperti ``memahami'' atau ``bernalar'', dan melupakan bahwa kecerdasan manusia tertanam dalam tubuh,
emosi, dan budaya, bukan hanya di otak. Mobil swakemudi dan robot rumah tangga yang dijanjikan sejak lama ternyata jauh lebih sulit
dari yang dibayangkan, dan pelajaran itu, menurutnya, berlaku juga untuk ramalan tentang AGI.

Kritik lain menyoroti bahwa perhatian pada risiko jangka jauh bisa mengalihkan perhatian dari kerugian yang sudah nyata hari ini:
bias terhadap kelompok yang kurang terwakili, pengawasan, pekerja yang memberi label data dengan upah rendah, dan biaya lingkungan
model besar \citep{bender2021}. Ada juga yang menunjuk bahwa sebagian narasi risiko ekstrem kebetulan menguntungkan perusahaan besar,
karena menggambarkan produk mereka sebagai kuat dan membenarkan regulasi yang menyulitkan pesaing kecil.

Buku ini tidak bisa memutuskan perdebatan itu, dan memang tidak perlu. Yang bisa dikatakan dengan cukup yakin adalah tiga hal. Pertama,
masalah-masalah teknis di balik kekhawatiran itu, seperti reward hacking, pergeseran distribusi, dan sulitnya memahami isi model, adalah
masalah nyata yang sudah terlihat pada sistem sekarang, dan dibahas di banyak bab buku ini. Kedua, ketidakpastian tentang masa depan
justru alasan untuk berhati-hati, bukan alasan untuk tenang. Ketiga, risiko masa kini dan risiko masa depan tidak harus bersaing:
banyak alat yang sama, seperti evaluasi yang ketat, transparansi, dan pengawasan manusia yang bermakna, membantu keduanya.

\section{Riset keselamatan yang konkret}

Di balik istilah-istilah besar, sebagian besar riset keselamatan AI berupa pekerjaan teknis yang biasa. Pemodelan preferensi dan
penyelarasan mencoba menangkap apa yang sebenarnya diinginkan manusia (\cref{ch:terbaru}). Interpretabilitas, termasuk sparse
autoencoder di \cref{ch:xai}, mencoba membaca apa yang terjadi di dalam model, sehingga perilaku berbahaya bisa dideteksi sebelum
muncul. Pengawasan yang bisa diskalakan bertanya bagaimana manusia bisa menilai jawaban yang terlalu panjang atau terlalu teknis untuk
diperiksa langsung, misalnya dengan bantuan model lain. Evaluasi kemampuan berbahaya menguji apakah sebuah model bisa, misalnya,
membantu membuat senjata biologis atau meretas sistem, sebelum model itu dirilis. \emph{Red-teaming} melibatkan orang yang sengaja
mencoba membuat model berperilaku buruk. Beberapa pengembang besar menerbitkan kerangka keselamatan yang menetapkan ambang kemampuan
dan tindakan pengaman yang harus dipenuhi sebelum model yang lebih kuat dilatih atau dirilis, dan AI Act mewajibkan evaluasi semacam
ini bagi model serbaguna dengan risiko sistemik (\cref{ch:manusia}).

Banyak dari pekerjaan ini memakai alat yang sudah kita kenal. Evaluasi yang jujur menuntut pembagian data tanpa kebocoran
(\cref{ch:data}). Pengukuran ketidakpastian membantu model tahu kapan ia tidak tahu (\cref{ch:stat}). Dan pelajaran dari keselamatan
sistem teknik lama, seperti sistem kendali pesawat atau reaktor, berlaku juga: sistem yang aman dirancang berlapis, dengan pengaman
yang tidak bergantung pada satu komponen yang sempurna.

\section{Tata kelola global}

Upaya internasional untuk mengatur AI tingkat lanjut bergerak cepat. Pada November 2023, Inggris menyelenggarakan AI Safety Summit
pertama di Bletchley Park, dihadiri wakil dari 28 negara serta Uni Eropa, termasuk Amerika Serikat dan Tiongkok. Pertemuan berikutnya
diadakan di Seoul pada Mei 2024, lalu di Paris pada Februari 2025 dengan nama AI Action Summit, dan di New Delhi pada Februari 2026
dengan nama AI Impact Summit, yang menghasilkan deklarasi yang didukung lebih dari 80 negara dan organisasi internasional. Pergantian
nama dari \emph{safety} ke \emph{action} lalu \emph{impact} mencerminkan pergeseran penekanan, dari risiko ke pemanfaatan dan
pemerataan, yang tidak semua pihak sambut dengan gembira.

Pertemuan di Bletchley juga mengamanatkan sebuah laporan ilmiah bersama. International AI Safety Report pertama, yang dipimpin Yoshua
Bengio dan disusun seratus ahli dengan wakil dari tiga puluh negara, PBB, OECD, dan Uni Eropa, terbit pada Januari 2025
\citep{bengio2025iasr}, dan edisi keduanya pada Februari 2026 \citep{iasr2026}. Laporan seperti ini meniru peran laporan IPCC untuk
iklim: bukan merekomendasikan kebijakan tertentu, tetapi merangkum apa yang diketahui dan apa yang belum, sebagai pijakan bersama bagi
pembuat kebijakan.

Satu perdebatan tata kelola yang terus berlanjut adalah model terbuka melawan model tertutup. Bobot model yang dirilis terbuka
memungkinkan riset, pemeriksaan independen, dan penggunaan di negara atau lembaga yang tidak mampu membayar layanan komersial, tetapi
juga tidak bisa ditarik kembali kalau ternyata disalahgunakan. Model tertutup lebih mudah dikendalikan, tetapi memusatkan kekuatan
pada segelintir perusahaan.

\section{Dampak yang sudah terasa}

Di luar perdebatan tentang masa depan, beberapa dampak sudah bisa diukur.

\emph{Pekerjaan.} \citet{eloundou2023} menilai pekerjaan di Amerika Serikat menurut seberapa cocok tugas-tugasnya dengan kemampuan model
bahasa besar. Mereka memperkirakan sekitar 80 persen pekerja punya setidaknya sepuluh persen tugas yang bisa terdampak, dan sekitar 19
persen punya setidaknya separuh tugas yang terdampak, dengan paparan yang justru lebih besar pada pekerjaan berpenghasilan tinggi.
Penulisnya menekankan bahwa ini ukuran paparan, bukan ramalan berapa pekerjaan yang hilang. Terdampak bisa berarti dibantu, diubah,
atau digantikan, dan sejarah teknologi sebelumnya memperlihatkan ketiganya terjadi bersamaan.

\emph{Energi.} Melatih dan menjalankan model besar membutuhkan listrik. \citet{strubell2019} termasuk yang pertama menghitung biaya
finansial dan jejak karbon pelatihan model NLP. Menurut perkiraan Badan Energi Internasional \citep{iea2025}, pusat data dunia memakai
sekitar 415 terawatt-jam listrik pada 2024, sekitar satu setengah persen konsumsi listrik global, dan dalam skenario dasarnya angka itu
naik menjadi sekitar 945 terawatt-jam pada 2030, dengan AI sebagai pendorong utama pertumbuhannya. Efisiensi model, seperti kuantisasi
dan distilasi di \cref{ch:terbaru} serta arsitektur dengan ingatan berukuran tetap di \cref{ch:lstm}, karenanya juga persoalan lingkungan.

\emph{Hak cipta dan data.} Model besar dilatih pada teks, gambar, dan kode yang sebagian besar punya pemilik. Sengketa hukum sudah
berjalan, misalnya gugatan The New York Times terhadap OpenAI dan Microsoft sejak akhir 2023, dan AI Act mewajibkan penyedia model
serbaguna punya kebijakan untuk menghormati hak cipta serta menerbitkan ringkasan data latihnya (\cref{ch:manusia}).

\emph{Informasi.} Model generatif membuat gambar, suara, dan video palsu yang meyakinkan menjadi murah. Watermark (\cref{ch:manusia}),
pelabelan konten sintetis yang diwajibkan regulasi, dan literasi AI adalah tiga lapis pertahanan yang saling melengkapi, dan tidak satu
pun cukup sendirian.

\emph{Pendidikan.} Siswa dan mahasiswa kini punya asisten yang bisa menulis esai dan menyelesaikan soal. Pertanyaan yang muncul bukan
lagi apakah alat ini dipakai, melainkan bagaimana belajar tetap terjadi ketika jawabannya tersedia seketika. Buku ini, dengan penurunan
rumus langkah demi langkah dan contoh hitungan, berangkat dari keyakinan sederhana bahwa memahami mengapa sebuah jawaban benar tetap
sesuatu yang harus dikerjakan sendiri.

\section{Penutup bab}

Perdebatan tentang AGI dan superintelligence sering terasa jauh dari pekerjaan sehari-hari menyetel jaringan atau membersihkan data.
Tetapi pertanyaan intinya sama dengan yang muncul di setiap bab buku ini: apakah model benar-benar mengerjakan apa yang kita maksud,
bagaimana kita tahu, dan apa yang terjadi di luar data yang pernah ia lihat. Semakin kuat sistemnya, semakin mahal jawaban yang salah.

\sumberbab{Bab ini disusun dari bacaan di luar kuliah, dengan sambungan ke kuliah AI and Law, Robopsychology, Reinforcement Learning, dan
AI and Visualization. Rujukan utama: \citet{bostrom2014}, \citet{russell2019}, \citet{amodei2016}, \citet{hendrycks2023},
\citet{mitchell2021}, \citet{morris2023}, laporan International AI Safety Report \citep{bengio2025iasr,iasr2026}, \citet{eloundou2023},
dan \citet{iea2025}. Tanggal dan angka peristiwa diperiksa dengan sumber resminya pada saat penulisan, Oktober 2026.}
